\documentclass[11pt,a4paper]{article}

\usepackage[utf8]{inputenc}
\usepackage[T1]{fontenc}
\usepackage[french]{babel}
\usepackage[a4paper,margin=2cm]{geometry}
\usepackage{amsmath,amssymb,amsthm,mathtools}
\usepackage{lmodern}
\usepackage{enumitem}
\usepackage{fancyhdr}
\usepackage{lastpage}
\usepackage{hyperref}
\usepackage{xcolor}
\usepackage{array}
\usepackage{tabularx}
\usepackage{multicol}

\setlength{\headheight}{14pt}
\pagestyle{fancy}
\fancyhf{}
\fancyhead[L]{Terminale générale -- Spécialité Mathématiques + Maths expertes}
\fancyhead[C]{Devoir surveillé -- Version corrigée}
\fancyhead[R]{Page \thepage/\pageref{LastPage}}
\fancyfoot[C]{\textit{La qualité du raisonnement, de la rédaction et de l'organisation de la preuve est essentielle}}

\newcommand{\R}{\mathbb{R}}
\newcommand{\C}{\mathbb{C}}
\newcommand{\Z}{\mathbb{Z}}
\newcommand{\N}{\mathbb{N}}
\newcommand{\e}{\mathrm{e}}
\newcommand{\ii}{\mathrm{i}}
\newcommand{\ds}{\displaystyle}

\begin{document}

\begin{center}
    {\Large \textbf{DEVOIR SURVEILLÉ -- Mathématiques}}\\[0.3em]
    {\large \textbf{Niveau Terminale générale + option Maths expertes}}\\[0.4em]
    {\large \textbf{Sujet d'approfondissement -- version corrigée}}\\[0.5em]
    Durée conseillée : 3h à 3h30
\end{center}

\vspace{0.4em}

\noindent\textbf{Consignes générales}
\begin{itemize}[leftmargin=1.8em]
    \item Tous les exercices sont indépendants.
    \item Les résultats doivent être justifiés avec soin.
    \item Les initiatives, les changements de point de vue et les raisonnements personnels pertinents sont valorisés.
    \item Une réponse juste mais non justifiée pourra être très peu valorisée.
\end{itemize}

\vspace{0.4em}

\begin{center}
\begin{tabular}{|>{\centering\arraybackslash}m{2.6cm}|>{\centering\arraybackslash}m{2.6cm}|>{\centering\arraybackslash}m{2.6cm}|>{\centering\arraybackslash}m{2.6cm}|}
\hline
Exercice 1 & Exercice 2 & Exercice 3 & Exercice 4 \\
\hline
18 points & 18 points & 22 points & 22 points \\
\hline
\end{tabular}
\end{center}

\vspace{0.5em}
\hrule
\vspace{0.6em}

\section*{Exercice 1 -- Fonctions, inégalités, raisonnement global \hfill /18}

On considère la fonction \(g\) définie sur \(]0,+\infty[\) par
\[
g(x)=x-\ln x.
\]

\subsection*{Partie A -- Étude fine d'une fonction}
\begin{enumerate}[label=\textbf{1.\arabic*}, leftmargin=1.6em]
    \item Déterminer les limites de \(g\) en \(0^+\) et en \(+\infty\).
    \item Étudier les variations de \(g\) sur \(]0,+\infty[\).
    \item Montrer que, pour tout réel \(x>0\),
    \[
    x-\ln x\ge 1.
    \]
    Préciser le cas d'égalité.
\end{enumerate}

\subsection*{Partie B -- Une inégalité moins immédiate}
On considère maintenant la fonction \(f\) définie sur \(]0,+\infty[\setminus\{1\}\) par
\[
f(x)=\frac{\ln x}{x-1}.
\]

\begin{enumerate}[label=\textbf{1.\arabic*}, resume, leftmargin=1.6em]
    \item Déterminer le signe de \(f(x)\) selon les valeurs de \(x\).
    \item Montrer que, pour tout \(x>0\), \(x\neq 1\),
    \[
    f'(x)=\frac{\frac{x-1}{x}-\ln x}{(x-1)^2}.
    \]
    \item En utilisant la fonction \(g\) ou toute autre méthode, étudier le signe de
    \[
    \frac{x-1}{x}-\ln x.
    \]
    \item En déduire les variations de \(f\) sur \(]0,1[\) puis sur \(]1,+\infty[\).
    \item Déterminer
    \[
    \lim_{x\to 1}\frac{\ln x}{x-1}.
    \]
    \item Montrer finalement que, pour tout réel \(x>0\),
    \[
    \ln x \le x-1.
    \]
    \item Étudier le cas d'égalité.
\end{enumerate}

\subsection*{Partie C -- Exploitation}
\begin{enumerate}[label=\textbf{1.\arabic*}, resume, leftmargin=1.6em]
    \item Montrer que, pour tout réel \(x>0\),
    \[
    \ln x \le \sqrt{x}-1
    \]
    ou infirmer cette assertion par un contre-exemple soigneusement justifié.
    \item Étudier le nombre de solutions de l'équation
    \[
    \ln x = \frac{x-1}{2}
    \]
    sur \(]0,+\infty[\).
\end{enumerate}

\vspace{0.5em}
\hrule
\vspace{0.6em}

\section*{Exercice 2 -- Probabilités, suite récurrente, vitesse de convergence \hfill /18}

Une plateforme en ligne propose aléatoirement une réduction à un utilisateur qui se connecte plusieurs jours de suite.

Pour tout entier naturel \(n\ge 1\), on note \(A_n\) l'événement :
\[
\text{\og l'utilisateur reçoit une réduction le \(n\)-ième jour \fg}
\]
et on pose
\[
a_n=P(A_n).
\]

On sait que :
\begin{itemize}
    \item au premier jour, la probabilité de recevoir une réduction est
    \[
    a_1=\frac12;
    \]
    \item si l'utilisateur a reçu une réduction un jour donné, alors la probabilité d'en recevoir une le lendemain est \(\dfrac13\) ;
    \item s'il n'en a pas reçu un jour donné, alors la probabilité d'en recevoir une le lendemain est \(\dfrac14\).
\end{itemize}

\subsection*{Partie A -- Mise en équation}
\begin{enumerate}[label=\textbf{2.\arabic*}, leftmargin=1.6em]
    \item Représenter la situation par un arbre entre les jours \(n\) et \(n+1\).
    \item Montrer que, pour tout entier \(n\ge 1\),
    \[
    a_{n+1}=\frac{1}{12}a_n+\frac14.
    \]
\end{enumerate}

\subsection*{Partie B -- Étude théorique}
\begin{enumerate}[label=\textbf{2.\arabic*}, resume, leftmargin=1.6em]
    \item Déterminer le réel \(\ell\) vérifiant
    \[
    \ell=\frac{1}{12}\ell+\frac14.
    \]
    \item On pose, pour tout \(n\ge 1\),
    \[
    u_n=a_n-\ell.
    \]
    Montrer que \((u_n)\) est géométrique.
    \item Préciser sa raison et son premier terme.
    \item En déduire l'expression de \(a_n\) en fonction de \(n\).
    \item Déterminer la limite de la suite \((a_n)\).
\end{enumerate}

\subsection*{Partie C -- Lecture qualitative}
\begin{enumerate}[label=\textbf{2.\arabic*}, resume, leftmargin=1.6em]
    \item Étudier le sens de variation de la suite \((a_n)\).
    \item Déterminer, selon les valeurs de \(n\), si \(a_n\) est supérieur ou inférieur à \(\ell\).
    \item Interpréter concrètement la valeur de \(\ell\).
\end{enumerate}

\subsection*{Partie D -- Question plus fine}
\begin{enumerate}[label=\textbf{2.\arabic*}, resume, leftmargin=1.6em]
    \item Déterminer le plus petit entier \(n\) tel que
    \[
    |a_n-\ell|<10^{-4}.
    \]
    \item On note
    \[
    S_n=a_1+a_2+\cdots+a_n.
    \]
    Donner une expression de \(S_n\) en fonction de \(n\).
    \item Interpréter \(S_n\) dans le contexte.
\end{enumerate}

\vspace{0.5em}
\hrule
\vspace{0.6em}

\section*{Exercice 3 -- Nombres complexes et géométrie : structure cachée \hfill /22}

Dans le plan complexe rapporté à un repère orthonormé direct, on considère les points
\(A\), \(B\), \(C\) d'affixes respectives
\[
z_A=1,\qquad z_B=\ii,\qquad z_C=-1.
\]

\subsection*{Partie A -- Première lecture géométrique}
\begin{enumerate}[label=\textbf{3.\arabic*}, leftmargin=1.6em]
    \item Représenter les points \(A\), \(B\) et \(C\).
    \item Déterminer les longueurs \(AB\), \(BC\) et \(CA\).
    \item En déduire la nature du triangle \(ABC\).
\end{enumerate}

\subsection*{Partie B -- Une transformation complexe}
On considère l'application \(T\) qui, à tout point \(M\) d'affixe \(z\), associe le point \(M'\) d'affixe
\[
z'=(1-\ii)z+1+\ii.
\]

\begin{enumerate}[label=\textbf{3.\arabic*}, resume, leftmargin=1.6em]
    \item Déterminer le module et un argument de \(1-\ii\).
    \item Interpréter géométriquement l'application
    \[
    z\mapsto (1-\ii)z.
    \]
    \item Montrer que l'application \(T\) admet un unique point fixe \(\Omega\) dont on déterminera l'affixe.
    \item Montrer que \(T\) est une similitude directe de centre \(\Omega\).
    \item Déterminer le rapport et l'angle de cette similitude.
\end{enumerate}

\subsection*{Partie C -- Itérations}
On définit une suite de points \((M_n)\) par :
\[
z_0=2,\qquad z_{n+1}=(1-\ii)z_n+1+\ii.
\]

\begin{enumerate}[label=\textbf{3.\arabic*}, resume, leftmargin=1.6em]
    \item Calculer \(z_1\) puis \(z_2\).
    \item On note \(\omega\) l'affixe du point fixe de \(T\). Montrer que, pour tout entier naturel \(n\),
    \[
    z_{n+1}-\omega=(1-\ii)(z_n-\omega).
    \]
    \item En déduire l'expression de \(z_n-\omega\) en fonction de \(n\).
    \item Montrer que
    \[
    |z_n-\omega|=(\sqrt2)^n|z_0-\omega|.
    \]
    \item Que peut-on en déduire sur le comportement géométrique des points \(M_n\) ?
\end{enumerate}

\subsection*{Partie D -- Une question de lieu}
On considère l'ensemble \(\mathcal E\) des points \(M\) d'affixe \(z\neq 1\) tels que
\[
\left|\frac{z-\ii}{z-1}\right|=\sqrt2.
\]

\begin{enumerate}[label=\textbf{3.\arabic*}, resume, leftmargin=1.6em]
    \item Donner une interprétation géométrique de cette relation.
    \item Montrer que \(\mathcal E\) est un cercle dont on précisera le centre et le rayon.
    \item Déterminer si le point \(C\) appartient à \(\mathcal E\).
\end{enumerate}

\vspace{0.5em}
\hrule
\vspace{0.6em}

\section*{Exercice 4 -- Arithmétique, congruences, structure des entiers \hfill /22}

\subsection*{Partie A -- Étude modulo 7}
\begin{enumerate}[label=\textbf{4.\arabic*}, leftmargin=1.6em]
    \item Calculer les restes modulo \(7\) des puissances successives de \(3\) :
    \[
    3,\ 3^2,\ 3^3,\ 3^4,\ 3^5,\ 3^6.
    \]
    \item En déduire que, pour tout entier naturel \(n\),
    \[
    3^6\equiv 1 \pmod 7.
    \]
    \item Déterminer tous les entiers naturels \(n\) tels que
    \[
    3^n\equiv 5 \pmod 7.
    \]
\end{enumerate}

\subsection*{Partie B -- Une équation plus structurée}
On cherche les couples \((x,y)\in\Z^2\) tels que
\[
3x-7y=5.
\]

\begin{enumerate}[label=\textbf{4.\arabic*}, resume, leftmargin=1.6em]
    \item Déterminer une solution particulière.
    \item Déterminer toutes les solutions dans \(\Z^2\).
    \item Déterminer toutes les solutions vérifiant de plus
    \[
    0\le x\le 20.
    \]
\end{enumerate}

\subsection*{Partie C -- Carrés modulo 8}
\begin{enumerate}[label=\textbf{4.\arabic*}, resume, leftmargin=1.6em]
    \item Montrer que, pour tout entier relatif \(n\), le reste de la division de \(n^2\) par \(8\) appartient à l'ensemble
    \[
    \{0,1,4\}.
    \]
    \item En déduire que l'équation
    \[
    n^2\equiv 3 \pmod 8
    \]
    n'admet aucune solution entière.
    \item Résoudre dans \(\Z\) la congruence
    \[
    n^2\equiv 1 \pmod 8.
    \]
\end{enumerate}

\subsection*{Partie D -- Problème de synthèse}
On considère l'équation
\[
x^2+7y^2=2026.
\]

\begin{enumerate}[label=\textbf{4.\arabic*}, resume, leftmargin=1.6em]
    \item Montrer que l’étude modulo \(8\) est pertinente pour ce problème.
    \item En étudiant l'équation modulo \(8\), montrer que l'équation n'admet aucune solution dans \(\Z^2\).
\end{enumerate}

\vspace{0.7em}
\hrule
\vspace{0.4em}

\section*{Barème indicatif}

\begin{itemize}[leftmargin=1.8em]
    \item Exercice 1 : 18 points
    \begin{itemize}
        \item Partie A : 5 points
        \item Partie B : 9 points
        \item Partie C : 4 points
    \end{itemize}
    \item Exercice 2 : 18 points
    \begin{itemize}
        \item Partie A : 3 points
        \item Partie B : 6 points
        \item Partie C : 4 points
        \item Partie D : 5 points
    \end{itemize}
    \item Exercice 3 : 22 points
    \begin{itemize}
        \item Partie A : 4 points
        \item Partie B : 7 points
        \item Partie C : 6 points
        \item Partie D : 5 points
    \end{itemize}
    \item Exercice 4 : 22 points
    \begin{itemize}
        \item Partie A : 5 points
        \item Partie B : 5 points
        \item Partie C : 5 points
        \item Partie D : 7 points
    \end{itemize}
\end{itemize}

\end{document}